\subsection{Assignment 18 -- Confronto stagionale delle categorie rispetto alla media dell’anno precedente}

L’Assignment 18 richiedeva di individuare, per ciascuna stagione, le categorie
musicali che nell’anno corrente registrano un numero di stream superiore alla
media della stessa stagione nell’anno precedente, riportando anche la relativa
variazione percentuale.

L’obiettivo dell’analisi è confrontare le performance stagionali delle singole
categorie con la tendenza complessiva del mercato nella medesima stagione
dell’anno precedente, così da evidenziare i generi che mostrano una crescita
significativa rispetto al benchmark stagionale.

Per costruire questa analisi, è stato innanzitutto definito un membro calcolato
che determina la media degli stream della stagione nell’anno precedente.
Tale valore è ottenuto utilizzando la funzione \textit{PARALLELPERIOD} per
risalire alla stessa stagione dell’anno precedente e la funzione \textit{AVG}
per calcolare la media degli stream su tutte le categorie. In questo contesto,
la “media della stagione” è stata interpretata come la media dell’intero mercato
stagionale e non come una media storica specifica della singola categoria.

Questo approccio consente di confrontare ogni categoria rispetto alla performance
complessiva del mercato nella stessa stagione, permettendo di identificare i
generi che performano meglio della tendenza stagionale generale.

Successivamente, è stata calcolata la variazione percentuale degli stream di
ciascuna categoria rispetto alla media stagionale dell’anno precedente. Il
calcolo include una gestione esplicita dei casi in cui il valore storico risulti
assente o pari a zero, al fine di evitare risultati non significativi o errori
di divisione.

Infine, è stato applicato un filtro per selezionare esclusivamente le categorie
che superano effettivamente la media della stagione precedente. Contestualmente,
sono state escluse la stagione \textit{Unknown} e le combinazioni prive di dati,
garantendo un output finale pulito e interpretabile dal punto di vista analitico.
